\section{Introduction}
\label{sec:intro}

The rise of digital social platforms has made \emph{opinion formation} a central
problem in the computational study of social systems. Individuals broadcast their
views, encounter those of others, and may revise their beliefs in response---a feedback
loop whose long-run outcome, broad consensus or entrenched polarization, depends in
subtle ways on the structure of the network and the behavioral rules governing each
interaction. Formal mathematical models of this process allow one to characterize
precisely the conditions under which each outcome arises.

The DeGroot model~\cite{DeGroot1974} is one of the most prominent formalisms
for social learning and opinion formation~\cite{Castellano2009}. A community is
represented as a weighted directed graph, the \emph{influence graph}, whose
nodes are \emph{agents} each holding an opinion in $[0,1]$---indicating the
strength of their agreement with a given proposition---and whose edges encode
how strongly agents influence one another. Agents repetitively revise their
beliefs by taking a weighted average of the opinions of those who influence them.
A fundamental theoretical result guarantees convergence to \emph{consensus}
whenever the influence graph is strongly connected and agents retain
non-zero self-influence~\cite{DeGroot1974,Castellano2009}. The significance
of this result lies in the fact that consensus is a central indicator of a
non-polarized community; indeed, the inability to reach consensus is a sign
of a polarized community.

The DeGroot model makes, however, two assumptions that can be overly constraining
in real social network contexts. First, it assumes that agents \emph{always}
express their opinions freely, disregarding the social pressure that may lead
individuals to withhold their views. Second, it assumes that agents
\emph{linearly} assimilate the opinions of those who influence them, ignoring
the nonlinear psychological responses that shape how people receive and react
to disagreement. Two well-documented phenomena challenge these assumptions,
respectively: the \emph{Spiral of Silence} and \emph{cognitive biases}.

The Spiral of Silence, introduced by Noelle-Neumann~\cite{NoelleNeumann1974,
NoelleNeumann1984}, posits that individuals whose opinions are in the
perceived minority tend to \emph{self-censor}: they withhold expression to
avoid social isolation, thereby reinforcing the dominance of the majority
view. Empirical studies confirm that this mechanism is active in both offline
and online contexts~\cite{Glynn1997,Hampton2014,Hampton2019,Neubaum2018,Sohn2022},
and simulation studies have shown that it can substantially distort the
collective opinion landscape in digitally mediated
networks~\cite{Sohn2022,Wang2013,Ross2019,Cabrera2021}. A formal treatment
was recently provided by Aranda et al.~\cite{aranda2024sound}, who extended the DeGroot
framework with silence mechanisms and derived conditions under which consensus
forms or fails in the presence of self-censorship. That work is the direct
predecessor of the present paper.

Cognitive biases introduce a further layer of nonlinearity into opinion
formation. When processing information that conflicts with their existing
beliefs, individuals do not respond neutrally: \emph{confirmation bias} leads
them to discount disagreeing views; the \emph{backfire effect} causes them to
harden their positions when confronted with contrary evidence; \emph{authority
bias} amplifies the influence of perceived high-status sources; and
\emph{insular bias} reduces the perceived relevance of distant
opinions~\cite{Aronson2010,Nyhan2010,Friedrich2020,Alvim2024}. These
distortions systematically alter how agents assimilate and reject information,
complicating consensus formation and fostering polarization in ways the
standard DeGroot framework cannot capture.

Alvim et al.~\cite{Alvim2024} introduced a formal framework that encodes these four biases
as nonlinear \emph{bias functions} applied to opinion disagreements, and
derived conditions under which consensus is preserved or fails in synchronous
DeGroot dynamics. However, that framework does not incorporate silence
mechanisms. Conversely, the model of Aranda et al.~\cite{aranda2024sound} addresses silence
but not cognitive biases. \textbf{The central contribution of the present
paper is to unify both dimensions}: we introduce two formal
Cognitive-Bias Silence Opinion Models (CBSOM)---\CBSOMminus\ (memoryless) and \CBSOMplus\ (memory-based)---that
embed cognitive bias functions into the Spiral-of-Silence opinion dynamics
of Aranda et al.~\cite{aranda2024sound}, and systematically analyze when the consensus
guarantees of the base model are preserved, altered, or broken under biased
processing and silence.

\noindent\textbf{Contributions.} The contributions of this paper are the
following:

\begin{enumerate}

  \item We introduce \CBSOMminus\ and \CBSOMplus, two Cognitive-Bias Silence
    Opinion Models that couple confirmation, backfire, authority, and insular
    bias functions with memoryless and memory-based belief update rules,
    respectively, within the Spiral-of-Silence--DeGroot framework
    (Section~\ref{sec:model}).

  \item We prove that the structural properties of the base Silence Opinion
    Model~\cite{aranda2024sound}---specifically, the impossibility of
    perpetual global silence and the monotonic contraction of opinion
    extremes---are preserved under all continuous bias functions in the
    \emph{receptive-resistant} region (Section~\ref{sec:analysis}).

  \item For the \emph{scaled linear} bias subfamily
    $\mathcal{F}_\alpha = \{f_\alpha(x)=\alpha x,\,\alpha\in(0,1)\}$,
    we prove consensus in cliques for all $n\ge 2$, including resolution
    of the period-2 oscillation pathology that affects dyadic cliques in
    the base model (Section~\ref{sec:analysis}).

  \item Through large-scale agent-based simulations on uniform and
    non-uniform networks with up to $10^{6}$ agents, we show that
    (i)~receptive-resistant biases (instantiated as $\mathrm{conf}(x)$ in our experiments) largely
    preserve consensus behavior, (ii)~backfire and authority biases
    robustly sustain polarization even in dense graphs, and
    (iii)~network size and topology interact with bias type in nuanced
    ways that are not predicted by small-scale analyses
    (Section~\ref{sec:experiments}).

\end{enumerate}

\noindent\textbf{Organization.}
Section~\ref{sec:model} introduces the theoretical background---the
DeGroot model, the cognitive bias framework, and the base Silence Opinion
Model of Aranda et al.~\cite{aranda2024sound}---and defines \CBSOMminus\ and \CBSOMplus.
Section~\ref{sec:analysis} presents the formal analysis.
Section~\ref{sec:experiments} reports the large-scale simulation study.
Section~\ref{sec:conclusions} concludes and discusses related work.
Complete proofs are collected in Appendix~\ref{app:proofs}.
